% J22 — Galois D_4 over LMFDB 4.2.10224.1: Number-Field Identification
%        of the Four-Core Attractor
%
% Brayden R. Sanders, M. Gish, 7Site LLC
% Submitted to: Communications in Algebra
% Source: extracted from J02 (four_core_consolidated.tex §8 — Galois D_4)
% Verification: sympy galois_group / sympy.factor(f, extension=[sqrt(-71)])
%               + LMFDB cross-check on 4.2.10224.1
% Companion cited as already-submitted: J02
% MSC 2020: 11R32, 11R16, 11R20, 12F10

\documentclass[11pt,reqno]{amsart}

\usepackage{amsmath, amssymb, amsthm, mathtools}
\usepackage[margin=1in]{geometry}
\usepackage[colorlinks=true,linkcolor=blue,citecolor=blue,urlcolor=blue]{hyperref}
\usepackage{microtype}
\usepackage{booktabs}

\theoremstyle{plain}
\newtheorem{theorem}{Theorem}[section]
\newtheorem{proposition}[theorem]{Proposition}
\newtheorem{lemma}[theorem]{Lemma}
\newtheorem{corollary}[theorem]{Corollary}

\theoremstyle{definition}
\newtheorem{definition}[theorem]{Definition}

\theoremstyle{remark}
\newtheorem*{remark}{Remark}

\newcommand{\Q}{\mathbb{Q}}
\newcommand{\Z}{\mathbb{Z}}
\newcommand{\Cfour}{\mathcal{C}_{4}}

\title[Galois $D_{4}$ over LMFDB 4.2.10224.1]
{Galois $D_{4}$ over LMFDB 4.2.10224.1: \\
 Number-Field Identification of the Four-Core Attractor}

\author{Brayden R.\ Sanders \and M.\ Gish}
\address{7Site LLC, Hot Springs, Arkansas, USA}
\email{brayden@7site.co}

\author{Brayden R.\ Sanders \and M.\ Gish}
\address{Independent Researcher, Hot Springs, Arkansas, USA}
\email{monica.gish1992@gmail.com}

\subjclass[2020]{11R32, 11R16, 11R20, 12F10}
\keywords{Galois group, dihedral extension, quartic number field,
LMFDB, fuse iteration, attractor identification}

\date{\today}

\begin{document}

\begin{abstract}
We isolate and present in self-contained form the Galois-theoretic
content of the four-core attractor identified in \cite{SandersGishFourCore}.
Specifically, the unique interior fixed point $p^{*}$ of the
symmetric-mixing fuse iteration $F_{1/2}$ on the four-core simplex
$\Delta^{3}_{\Cfour}$ has its coordinate ratio $\xi^{*} := r/\beta$
characterized as the unique positive real root of the irreducible monic
integer quartic
\[
  f(x) \;=\; x^{4} + 4x^{3} - x^{2} + 2x - 2.
\]
The Galois group of $f$ over $\Q$ is the dihedral group $D_{4}$ of order
$8$. The number field $K = \Q[x]/(f)$ has discriminant
$d_{K} = -2^{4}\cdot 3^{2}\cdot 71 = -10224$, class number
$h_{K} = 1$, signature $(2,1)$, regulator $R_{K} \approx 8.617$, and is
the catalogued field LMFDB~$4.2.10224.1$. The Galois closure
$L \supset K$ has degree $8$ over $\Q$, signature $(0,4)$, class number
$14$, discriminant $2^{8}\cdot 3^{4}\cdot 71^{4}$, and is
LMFDB~$8.0.526936617216.1$. The number field is catalogued and not new;
what is new is the route: a quartic Galois $D_{4}$ field arising as the
ring of definition for the fixed-point coordinates of a symmetric-mixing
iteration on a four-element fusion-closed sub-magma of $\Z/10$.
\end{abstract}

\maketitle

\tableofcontents

%==============================================================
\section{Introduction and statement of result}
%==============================================================

The four-core sub-magma $\Cfour = \{0, 7, 8, 9\} \subset \Z/10$ is
fusion-closed under the joint action of the two canonical commutative
binary operations $\TS, \BH$ on $\Z/10$ (the magma data and joint
closure are established in \cite{SandersGishFourCore}). The
symmetric-mixing fuse iteration $F_{1/2}$ on the four-core simplex
$\Delta^{3}_{\Cfour} = \{(v, h, \beta, r) \in \R_{\geq 0}^{4} :
v + h + \beta + r = 1\}$ has a unique interior fixed point $p^{*}$.
The Galois content of $p^{*}$ is captured by two algebraic ratios:
$h/\beta = 1 + \sqrt{3}$ and $\xi^{*} = r/\beta$.

The first ratio places one quadratic extension $\Q(\sqrt{3}) \subset
\Q(\xi^{*})$ in the Galois closure. The second ratio is the focus of
the present paper.

\begin{theorem}[Main result, restated from \cite{SandersGishFourCore} §8]
\label{thm:galois}
Let $p^{*}$ be the unique interior fixed point of the
symmetric-mixing fuse iteration $F_{1/2}$ on the four-core simplex
$\Delta^{3}_{\Cfour}$. Then the coordinate ratio
$\xi^{*} = p^{*}_{r}/p^{*}_{\beta}$ is the unique positive real root of
the irreducible monic integer quartic
\[
  f(x) \;=\; x^{4} + 4x^{3} - x^{2} + 2x - 2.
\]
The Galois group of $f$ over $\Q$ is the dihedral group $D_{4}$ of
order $8$. The number field $K = \Q[x]/(f)$ has field discriminant
$d_{K} = -2^{4}\cdot 3^{2}\cdot 71$, class number $h_{K} = 1$,
signature $(2, 1)$, regulator $R_{K} \approx 8.617$, and is the
catalogued field LMFDB~$4.2.10224.1$. The Galois closure $L \supset K$
has degree $8$ over $\Q$, signature $(0, 4)$, class number $14$,
discriminant $2^{8}\cdot 3^{4}\cdot 71^{4}$, and is
LMFDB~$8.0.526936617216.1$. The intermediate field $\Q(\sqrt{3}) \subset
L$ is a genuine subfield of $L$, with explicit factorization
\[
  f(x) \;=\; \bigl(x^{2}+(2-\sqrt{3})x+(\sqrt{3}-1)\bigr)
             \bigl(x^{2}+(2+\sqrt{3})x-(\sqrt{3}+1)\bigr)
\]
in $\Q(\sqrt{3})[x]$.
\end{theorem}

\subsection{Scope of the present paper.}
The four-core paper \cite{SandersGishFourCore} establishes
Theorem~\ref{thm:galois} as part of a wider analysis (joint closure of
$\TS, \BH$ on the four-core; per-coordinate fuse data; closed-form
attractor at $\alpha = 1/2$; $\alpha$-uniqueness on the
Stern-Brocot grid). The present paper extracts the Galois-theoretic
content as a self-contained result: definition of the iteration,
derivation of the quartic, irreducibility, Galois-group computation,
and number-field identification with LMFDB. We assume only the joint
fuse data of \cite{SandersGishFourCore}, which we recall in
Section~\ref{sec:setup}; the bulk of the present paper (Sections
\ref{sec:derivation}--\ref{sec:lmfdb}) reproduces the derivation in
self-contained form.

%==============================================================
\section{Setup: the four-core fuse iteration}
\label{sec:setup}
%==============================================================

We recall the structural input from \cite{SandersGishFourCore}.

\begin{definition}\label{def:fourcore-iteration}
Let $V = e_{0}, H = e_{7}, B = e_{8}, R = e_{9}$ be the four basis
vectors corresponding to the indices in $\Cfour = \{0, 7, 8, 9\}$.
A point in the four-core simplex is $p = (v, h, \beta, r) \in
\R_{\geq 0}^{4}$ with $v + h + \beta + r = 1$.
The (symmetric-mixing) fuse iteration $F_{1/2}: \Delta^{3}_{\Cfour}
\to \Delta^{3}_{\Cfour}$ is given by
\[
  F_{1/2}(p) \;=\; \tfrac12\bigl(\TS_{\rm fuse}(p) + \BH_{\rm fuse}(p)\bigr),
\]
where the two fuse operators are the per-coordinate fuse maps tabulated
in \cite{SandersGishFourCore} Proposition (per-coordinate fuse data).
\end{definition}

\begin{remark}
The crucial structural input we use below is the vanishing
$\TS_{\rm fuse}(p)_{9} = 0$ on the four-core simplex (the four-core is
fusion-closed under $\TS$ but $r$-coordinate inputs do not appear in the
$r$-coordinate output of $\TS_{\rm fuse}$). This vanishing, together
with the $\beta$-coordinate identity
$\BH_{\rm fuse}(p)_{9} = 2vr + 2h\beta$, is what produces the quadratic
relation in $\xi$.
\end{remark}

%==============================================================
\section{Derivation of the quartic}
\label{sec:derivation}
%==============================================================

\subsection{Derivation through $\Q(\sqrt{3})$.}

The fixed-point equation for the $r$-coordinate at $\alpha = \tfrac12$
gives, using the per-coordinate fuse data of
\cite{SandersGishFourCore} and the vanishing
$\TS_{\rm fuse}(p^{*})_{9} = 0$,
\[
  r \;=\; \tfrac12\bigl(\TS_{\rm fuse}(p^{*})_{9} + \BH_{\rm fuse}(p^{*})_{9}\bigr)
        \;=\; \tfrac12(2vr + 2 h\beta)
        \;=\; vr + h\beta.
\]
Substituting $v = 1 - h - \beta - r$ and rearranging,
\[
  r^{2} + (h + \beta)\, r \;-\; h\beta \;=\; 0.
\]
The closed-form $h/\beta = 1 + \sqrt{3}$ from \cite{SandersGishFourCore}
gives $h = (1+\sqrt{3})\beta$. Setting $\xi = r/\beta$ and dividing by
$\beta^{2}$,
\[
  \xi^{2} + (2+\sqrt{3})\,\xi \;-\; (1+\sqrt{3}) \;=\; 0
  \qquad\text{over }\Q(\sqrt{3}).
\]
The positive root $\xi^{*} > 0$ is therefore an algebraic number of
degree $2$ over $\Q(\sqrt{3})$.

\subsection{Norm to $\Q$.}

The minimal polynomial of $\xi^{*}$ over $\Q$ is the product of the
displayed quadratic with its $\mathrm{Gal}(\Q(\sqrt{3})/\Q)$-conjugate
(replacing $\sqrt{3}$ with $-\sqrt{3}$):
\begin{align*}
  f(x)
  &= \bigl(x^{2} + (2+\sqrt{3})x - (1+\sqrt{3})\bigr)
     \bigl(x^{2} + (2-\sqrt{3})x - (1-\sqrt{3})\bigr) \\
  &= x^{4} + 4x^{3} - x^{2} + 2x - 2.
\end{align*}
The factorization is the $\Q(\sqrt{3})$-factorization of $f$,
exhibited as part of the derivation rather than verified post-hoc.

%==============================================================
\section{Irreducibility over $\Q$}
\label{sec:irreducibility}
%==============================================================

By Gauss's lemma, any factorization of $f$ in $\Q[x]$ gives a
factorization in $\Z[x]$. Direct evaluation rules out integer roots:
$f(\pm 1) \in \{4, -8\}$ and $f(\pm 2) \in \{46, -26\}$, none zero.
For an integer factorization
$f = (x^{2} + ax + b)(x^{2} + cx + d)$, matching coefficients gives
\[
  a + c = 4,\quad ac + b + d = -1,\quad ad + bc = 2,\quad bd = -2.
\]
The constraint $bd = -2$ over $\Z$ has solutions
$(b, d) \in \{(1, -2), (-1, 2), (2, -1), (-2, 1)\}$. In each case the
linear equation $ad + bc = 2$ together with $c = 4 - a$ determines $a$,
and the quadratic $ac + b + d = -1$ gives a consistency check:
\begin{itemize}\setlength\itemsep{0pt}
\item $(b, d) = (1, -2)$: $-2a + (4 - a) = 2$ gives $a = 2/3 \notin \Z$.
\item $(b, d) = (-1, 2)$: $2a - (4 - a) = 2$ gives $a = 2 \in \Z$,
$c = 2$, but then $ac = 4 \neq -2 = -1 - b - d$, contradicting the
quadratic constraint.
\item $(b, d) = (2, -1)$: $-a + 2(4 - a) = 2$ gives $a = 2 \in \Z$,
$c = 2$, again $ac = 4 \neq -2$.
\item $(b, d) = (-2, 1)$: $a - 2(4 - a) = 2$ gives $a = 10/3 \notin \Z$.
\end{itemize}
No integer factorization exists; $f$ is irreducible over $\Z$, hence
over $\Q$.

\begin{remark}[Reduction-modulo-$p$ alternative]
$f$ is irreducible mod $7$ (one-line proof). Note that $f$ is
\emph{reducible} mod $5$ as
$f \equiv (x^{2} - 2)(x^{2} - x + 1) \pmod{5}$, so reduction-based
arguments require care in the choice of prime.
\end{remark}

%==============================================================
\section{Galois group computation}
\label{sec:galois}
%==============================================================

The polynomial discriminant is
\[
  \Delta_{f} \;=\; -2^{6}\cdot 3^{2}\cdot 71
\]
(computed directly). The resolvent cubic of $f$ is
\[
  g(y) \;=\; y^{3} + y^{2} + 16y + 36
       \;=\; (y + 2)(y^{2} - y + 18),
\]
with the unique rational root $y = -2$ and an irreducible quadratic
factor of discriminant $-71$.

By the standard quartic-Galois-group classification (Cohen,
\textit{A Course in Computational Algebraic Number Theory},
\S6.3.2~\cite{Cohen1993}), the existence of exactly one rational root
in the resolvent cubic places
$\mathrm{Gal}(f/\Q)\in\{C_{4}, D_{4}\}$. The two are distinguished by
whether $f$ remains irreducible over $\Q(\sqrt{\Delta_{f}}) =
\Q(\sqrt{-71})$: irreducibility there gives $D_{4}$, factorization
gives $C_{4}$. Direct computation
(\texttt{sympy.factor(f, extension=[sqrt(-71)])}) returns $f$ itself as
a single irreducible factor, confirming
\[
  \mathrm{Gal}(f/\Q) \;=\; D_{4}.
\]

\begin{remark}[Visibility of $D_{4}$ from the derivation]
The $D_{4}$-Galois group is also visible directly from the derivation
above. The closed form $h/\beta = 1+\sqrt{3}$ generates the quadratic
subfield $\Q(\sqrt{3})\subset\Q(\xi^{*})$, and the second nontrivial
quadratic subfield $\Q(\sqrt{-71})$ enters through the discriminant;
the splitting field is the compositum
$\Q(\sqrt{3}, \sqrt{-71}, \xi^{*})$, a $D_{4}$-extension of $\Q$.
\end{remark}

%==============================================================
\section{Number-field identification with LMFDB}
\label{sec:lmfdb}
%==============================================================

The field $K = \Q[x]/(f)$ has degree $4$ and signature $(2, 1)$
(counting real roots of $f$: two real, one complex conjugate pair).
The Tschirnhaus substitution $x \mapsto -x - 1$ converts $f$ to
\[
  h(x) \;=\; x^{4} - 7x^{2} - 12x - 8,
\]
the canonical defining polynomial recorded in LMFDB's
catalogue~\cite{LMFDB-4-2-10224-1} for the quartic field with
$d_{K} = -10224 = -2^{4}\cdot 3^{2}\cdot 71$, class number $h_{K} = 1$,
regulator $R_{K} \approx 8.617$, and Galois group $D_{4}$. Hence
\[
  K \;=\; \text{LMFDB}~4.2.10224.1.
\]
The Galois closure is the $D_{4}$-extension recorded
in~\cite{LMFDB-8-0-526936617216-1}; intermediate subfields include
$\Q(\sqrt{3})$, $\Q(\sqrt{-71})$, and the biquadratic
$\Q(\sqrt{3}, \sqrt{-71})$.

\begin{remark}[Novelty of the route, not the field]
The number field $K = $ LMFDB~$4.2.10224.1$ is catalogued and not new;
what is novel is the derivation. The polynomial form
$f(x) = x^{4} + 4x^{3} - x^{2} + 2x - 2$ does not appear in OEIS for
its coefficient sequence $[1, 4, -1, 2, -2]$ as a distinguished
sequence, and we have not located it as a defining polynomial in a
published context tied to a finite-magma fuse iteration. The pitch is
therefore: a new route to a known number field, with the ring of
definition realized as the fixed-point coordinates of the
symmetric-mixing-weight iteration $F_{1/2}$ on the four-core simplex
$\Delta^{3}_{\Cfour}$.
\end{remark}

%==============================================================
\section{Verification}
\label{sec:verification}
%==============================================================

The computational verification consists of:
\begin{enumerate}
\item Polynomial extraction from the fuse data
  (\texttt{07\_full\_closed\_form.py}, \cite{SandersGishFourCore}).
\item Verification of $D_{4}$ Galois group via
  \texttt{sympy.factor(f, extension=[sqrt(-71)])}.
\item Explicit factorization in $\Q(\sqrt{3})[x]$ via
  \texttt{sympy.factor(f, extension=[sqrt(3)])}.
\item LMFDB cross-check by Tschirnhaus reduction to LMFDB's defining
  polynomial $x^{4} - 7x^{2} - 12x - 8$.
\end{enumerate}

All four checks complete deterministically in $<5$ seconds. The
verification scripts and the derivation tables are deposited at
\url{https://github.com/TiredofSleep/ck/tree/tig-synthesis} alongside
the four-core consolidated paper \cite{SandersGishFourCore}.

%==============================================================
\section*{Acknowledgments}
%==============================================================

This work extracts and presents in self-contained form the Galois
content of the four-core attractor identified in
\cite{SandersGishFourCore}. The present manuscript is self-contained as
a Galois-theoretic study of the quartic
$x^{4} + 4x^{3} - x^{2} + 2x - 2$; familiarity with the broader
fuse-iteration framework is not required beyond
Section~\ref{sec:setup}. We thank the LMFDB collaboration for the
catalogue infrastructure that makes number-field identification
routine. Algorithmic verification was developed with assistance from
Anthropic Claude sessions in 2026.

%==============================================================
\begin{thebibliography}{99}
%==============================================================

\bibitem{SandersGishFourCore} B.~R.~Sanders, M.~Gish,
\emph{Joint Closure, Per-Coordinate Fuse Data, and a Closed-Form
Algebraic Attractor of Two Commutative Binary Operations on
$\Z/10$}, submitted to \emph{Algebraic Combinatorics}, 2026 (J02).

\bibitem{Cohen1993} H.~Cohen,
\emph{A Course in Computational Algebraic Number Theory},
Graduate Texts in Mathematics 138, Springer, 1993.

\bibitem{LMFDB-4-2-10224-1} The {LMFDB} Collaboration,
\emph{Number field 4.2.10224.1},
The L-functions and Modular Forms Database,
\url{https://www.lmfdb.org/NumberField/4.2.10224.1}.

\bibitem{LMFDB-8-0-526936617216-1} The {LMFDB} Collaboration,
\emph{Number field 8.0.526936617216.1},
The L-functions and Modular Forms Database,
\url{https://www.lmfdb.org/NumberField/8.0.526936617216.1}.

\bibitem{TIGsynthesis2026} B.~R.~Sanders,
\emph{Trinity Infinity Geometry Synthesis},
github.com/TiredofSleep/ck (default branch),
DOI: 10.5281/zenodo.18852047, 2026.

\end{thebibliography}

\end{document}
